\appendix

\chapter{Results}\label{app:results}

This appendix collects the detailed results that support Chapter~\ref{evaluation}.

\section{Per-Class Performance}

Table~\ref{tab:per-class} gives precision, recall, F1 and support for all 30
classes on the canonical test fold, for the prototype model
(\texttt{comparison/\allowbreak protgnn\_canonical}), sorted by support. It is
the detail behind the aggregate numbers in Table~\ref{tab:main-results} and the
error analysis in Section~\ref{sec:error-analysis}.

\begin{table}[htbp]
  \centering
  \small
  \setlength{\tabcolsep}{4pt}
  \caption{Per-class test performance of the prototype model on the canonical
    split. Macro-F1 is the unweighted mean of the F1 column (0.508).}
  \label{tab:per-class}
  \begin{tabular}{lrrrr}
    \toprule
    Diagnosis class & Precision & Recall & F1 & Support \\
    \midrule
    Limb injury or pain & 0.732 & 0.615 & 0.668 & 903 \\
    Back or spine pain & 0.820 & 0.644 & 0.722 & 855 \\
    Cardiovascular risk factor & 0.586 & 0.406 & 0.479 & 749 \\
    Head injury & 0.687 & 0.490 & 0.572 & 502 \\
    Diabetes mellitus & 0.584 & 0.551 & 0.567 & 425 \\
    UTI or pyelonephritis & 0.590 & 0.341 & 0.433 & 413 \\
    Skin or soft-tissue infection & 0.678 & 0.524 & 0.591 & 389 \\
    Lower respiratory disease & 0.464 & 0.366 & 0.409 & 369 \\
    Alcohol abuse with intoxication & 0.897 & 0.812 & 0.852 & 345 \\
    GI bleed & 0.810 & 0.666 & 0.731 & 302 \\
    Major depressive disorder & 0.832 & 0.835 & 0.834 & 267 \\
    Unspecified atrial fibrillation & 0.420 & 0.497 & 0.455 & 179 \\
    Acute pharyngitis & 0.828 & 0.744 & 0.784 & 168 \\
    Anxiety disorder & 0.497 & 0.601 & 0.544 & 138 \\
    Dehydration & 0.189 & 0.419 & 0.261 & 129 \\
    Acute kidney failure & 0.377 & 0.359 & 0.368 & 128 \\
    Hypokalemia & 0.266 & 0.504 & 0.348 & 127 \\
    Laceration w/o fb of l idx fngr w/o damage to nail & 0.449 & 0.754 & 0.563 & 122 \\
    Unsp intestnl obst & 0.591 & 0.556 & 0.573 & 117 \\
    Anemia & 0.325 & 0.693 & 0.443 & 114 \\
    Unspecified asthma with (acute) exacerbation & 0.481 & 0.670 & 0.560 & 112 \\
    Epilepsy & 0.608 & 0.784 & 0.685 & 97 \\
    Heart failure & 0.353 & 0.624 & 0.451 & 85 \\
    Acute upper respiratory infection & 0.175 & 0.488 & 0.257 & 82 \\
    Non-ST elevation (NSTEMI) myocardial infarction & 0.378 & 0.750 & 0.502 & 72 \\
    Multiple fractures of ribs & 0.223 & 0.529 & 0.314 & 70 \\
    Hypothyroidism & 0.236 & 0.600 & 0.339 & 65 \\
    Hypotension & 0.200 & 0.277 & 0.232 & 47 \\
    End stage renal disease & 0.404 & 0.477 & 0.438 & 44 \\
    Sepsis & 0.196 & 0.488 & 0.280 & 41 \\
    \bottomrule
  \end{tabular}
\end{table}

\section{Class Distribution}

The per-class counts for the canonical run are listed in
Table~\ref{tab:class-dist}, sorted by training frequency. The class names are
the merged disease categories used by the canonical split
(\texttt{comparison/\allowbreak canonical\_split.json}); earlier runs used the raw
\texttt{disease\_1} strings and are not comparable.

\begin{table}[htbp]
  \centering
  \small
  \caption{Class distribution for the 30-class disease dataset on the canonical
    split (train / validation / test).}
  \label{tab:class-dist}
  \begin{tabular}{lrrrr}
    \toprule
    Diagnosis class & Train & Share & Val & Test \\
    \midrule
    Limb injury or pain & 7{,}224 & 12.1\% & 903 & 903 \\
    Back or spine pain & 6{,}841 & 11.5\% & 855 & 855 \\
    Cardiovascular risk factor & 5{,}986 & 10.0\% & 748 & 749 \\
    Head injury & 4{,}016 & 6.7\% & 502 & 502 \\
    Diabetes mellitus & 3{,}395 & 5.7\% & 424 & 425 \\
    UTI or pyelonephritis & 3{,}298 & 5.5\% & 412 & 413 \\
    Skin or soft-tissue infection & 3{,}114 & 5.2\% & 389 & 389 \\
    Lower respiratory disease & 2{,}952 & 5.0\% & 369 & 369 \\
    Alcohol abuse with intoxication & 2{,}755 & 4.6\% & 344 & 345 \\
    GI bleed & 2{,}417 & 4.1\% & 302 & 302 \\
    Major depressive disorder & 2{,}141 & 3.6\% & 268 & 267 \\
    Unspecified atrial fibrillation & 1{,}428 & 2.4\% & 178 & 179 \\
    Acute pharyngitis & 1{,}346 & 2.3\% & 168 & 168 \\
    Anxiety disorder & 1{,}103 & 1.9\% & 138 & 138 \\
    Dehydration & 1{,}027 & 1.7\% & 128 & 129 \\
    Acute kidney failure & 1{,}025 & 1.7\% & 128 & 128 \\
    Hypokalemia & 1{,}012 & 1.7\% & 126 & 127 \\
    Laceration w/o fb of l idx fngr w/o damage to nail & 971 & 1.6\% & 121 & 122 \\
    Unsp intestnl obst & 940 & 1.6\% & 118 & 117 \\
    Anemia & 913 & 1.5\% & 114 & 114 \\
    Unspecified asthma with (acute) exacerbation & 894 & 1.5\% & 112 & 112 \\
    Epilepsy & 774 & 1.3\% & 97 & 97 \\
    Heart failure & 680 & 1.1\% & 85 & 85 \\
    Acute upper respiratory infection & 656 & 1.1\% & 82 & 82 \\
    Non-ST elevation (NSTEMI) myocardial infarction & 577 & 1.0\% & 72 & 72 \\
    Multiple fractures of ribs & 558 & 0.9\% & 70 & 70 \\
    Hypothyroidism & 515 & 0.9\% & 64 & 65 \\
    Hypotension & 381 & 0.6\% & 48 & 47 \\
    End stage renal disease & 346 & 0.6\% & 43 & 44 \\
    Sepsis & 322 & 0.5\% & 40 & 41 \\
    \midrule
    Total & 59{,}607 & 100.0\% & 7{,}448 & 7{,}456 \\
    \bottomrule
  \end{tabular}
\end{table}

\section{Example Explanations}

The three cases below come from the explanation pass over the canonical test
fold (\texttt{comparison/\allowbreak protgnn\_canonical/\allowbreak
graphxai\_summary.csv}). \caveat{That pass covers 50 graphs, 0.7\% of the test
fold, so these cases illustrate the mechanisms described in
Section~\ref{sec:expl-quality} rather than establishing how often each one
occurs.} They were chosen to show one correct prediction with coherent evidence,
one failure where the model reads the wrong node, and the relation between
prototype distance and correctness. Node labels are the decoded
clinical tokens: \texttt{cc[\ldots]} is a triage chief complaint,
\texttt{sym[\ldots]} a doctor-assigned symptom code, drug names are home or
ED-administered medications, and arrows give the z-scored direction of a vital
sign.

\paragraph{Case 1, correct prediction with coherent evidence (graph 2).}
Predicted and true class: \emph{Laceration w/o fb of l idx fngr w/o damage to
nail}; 17 nodes. Integrated Gradients ranks \texttt{cc[laceration]} first and the
gradient explainer ranks it second, behind the patient node. The three nearest
prototypes (\#56, \#55, \#54, distances 0.074 to 0.076) all belong to the
predicted class, so the case-based evidence and the attribution point the same
way. GNNExplainer, as on most graphs, ranks the demographic hub node first and
then vital signs, which says little about the diagnosis.

\paragraph{Case 2, the model reads the wrong node (graph 5).}
Predicted \emph{Hypokalemia}, true class \emph{Lower respiratory disease}; 22
nodes. The visit carries an oxygen saturation 4.8 standard deviations below the
mean, flagged abnormal, which is the signal for the true class. The gradient
explainer does surface it, but only after \texttt{cc[rash]} and the patient node;
Integrated Gradients drops it entirely in favour of levetiracetam and
lamotrigine, the patient's antiepileptic medication. The nearest prototype is
\#46 (\emph{Hypokalemia}) at distance 0.173, and the next two are also
hypokalaemia prototypes, so the model is confidently matching the wrong
archetype. The explanation is faithful, it shows what the model used, but what
the model used was the medication history rather than the vital sign.

\paragraph{Case 3, prototype distance as a confidence signal (graph 0).}
Predicted and true class: \emph{GI bleed}; 17 nodes. The nearest prototype is
\#38 at distance 0.032, the closest match in the sample, with the next two
prototypes of the same class at 0.054 and 0.058. Integrated Gradients selects
\texttt{cc[rectal pain]} first, which is the presenting complaint for this
diagnosis. This is the pattern behind the aggregate numbers in
Section~\ref{sec:expl-quality}: tight prototype distances (mean 0.122 \caveat{over
the 50 explained graphs}) go with correct predictions, while the loose matches (mean
0.191) are where the model fails, as in Case 2.
